%! Logarithmic Function Manipulation — Unit 5, Lesson 5.4 — Homework
\documentclass[10pt]{article}
\usepackage{precalculus-article}
\usepackage{precalculus-boxes}
\usepackage{grading}
\gradingassignment{precalc-u05-l04-hw}
\gradingcourse{precalculus}
\newcommand{\TallMath}[1]{$\displaystyle #1\rule[-1.4em]{0pt}{3.2em}$}

\begin{document}

\pageheader{Unit 5, Lesson 5.4}{Homework}

\vspace{0.1in}

\begin{enumerate}[leftmargin=1.8em, itemsep=14pt]

\item Apply the product property to rewrite each expression as a sum.

\begin{enumerate}[leftmargin=1.5em, itemsep=6pt, label=\alph*.]
  \item $\log_6(36 \cdot 216) = \log_6(\smash{\gradedblank[id=hw1a1, type=short, points=1, answer={36}]{2cm}{13pt}}) + \log_6(\smash{\gradedblank[id=hw1a2, type=short, points=1, answer={216}]{2cm}{13pt}}) = \smash{\gradedblank[id=hw1a3, type=short, points=1, answer={2}]{1cm}{13pt}} + \smash{\gradedblank[id=hw1a4, type=short, points=1, answer={3}]{1cm}{13pt}} = \smash{\gradedblank[id=hw1a5, type=short, points=1, answer={5}]{1.5cm}{13pt}}$

  \item $\log_2(x \cdot 8) = \log_2(x) + \log_2(\smash{\gradedblank[id=hw1b1, type=short, points=1, answer={8}]{2cm}{13pt}}) = \log_2(x) + \smash{\gradedblank[id=hw1b2, type=short, points=1, answer={3}]{1.5cm}{13pt}}$
\end{enumerate}

\item Apply the power property to rewrite each expression.

\begin{enumerate}[leftmargin=1.5em, itemsep=6pt, label=\alph*.]
  \item $\log_5(125^3) = \smash{\gradedblank[id=hw2a1, type=short, points=1, answer={3}]{1cm}{13pt}} \cdot \log_5(\smash{\gradedblank[id=hw2a2, type=short, points=1, answer={125}]{2cm}{13pt}}) = \smash{\gradedblank[id=hw2a3, type=short, points=1, answer={3}]{1cm}{13pt}} \cdot \smash{\gradedblank[id=hw2a4, type=short, points=1, answer={3}]{1cm}{13pt}} = \smash{\gradedblank[id=hw2a5, type=short, points=1, answer={9}]{1.5cm}{13pt}}$

  \item $\log_4(x^5) = \smash{\gradedblank[id=hw2b, type=short, points=1, answer={5\log_4(x)}]{1.5cm}{13pt}}$
\end{enumerate}

\item Let $f(x) = \log_2(4x)$.

\begin{enumerate}[leftmargin=1.5em, itemsep=6pt, label=\alph*.]
  \item Rewrite $f(x)$ in the form $f(x) = a + \log_2(x)$. What is $a$?

  \vspace{4pt}
  $f(x) = \log_2(4x) = \log_2(\smash{\gradedblank[id=hw3a1, type=short, points=1, answer={4}]{2cm}{13pt}}) + \log_2(x) = \smash{\gradedblank[id=hw3a2, type=short, points=1, answer={2}]{1.5cm}{13pt}} + \log_2(x)$

  \vspace{2pt}
  $a = $ \smash{\gradedblank[id=hw3a3, type=short, points=1, answer={2}]{2cm}{13pt}}

  \item What geometric transformation does this represent on the graph of $y = \log_2(x)$?

        \rubricbox{\gradedblank[id=hw3b, type=rubric, points=1,
  display={A vertical translation up 2 units of $y = \log_2(x)$.}]{\linewidth}{13pt}}
\end{enumerate}

\item Use the change-of-base formula to show that $\log_3(x) = \dfrac{\log_{10}(x)}{\log_{10}(3)}$.
      Explain why this means $\log_3(x)$ is a vertical dilation of $\log_{10}(x)$.

      \rubricbox{\gradedblank[id=hw4, type=rubric, points=2,
  display={Change of base: $\log_3(x) = \dfrac{\log_{10}(x)}{\log_{10}(3)} = \dfrac{1}{\log_{10}(3)}\cdot\log_{10}(x)$. Multiplying every output of $\log_{10}(x)$ by the constant $1/\log_{10}(3) \approx 2.096$ is a vertical dilation.}]{\linewidth}{40pt}}

\newpage
\item Evaluate each natural logarithm. Show reasoning.

\begin{enumerate}[leftmargin=1.5em, itemsep=6pt, label=\alph*.]
  \item $\ln(e^5) = $ \smash{\gradedblank[id=hw5a, type=short, points=1, answer={5}]{3cm}{13pt}}

  \item $\ln(1) = $ \smash{\gradedblank[id=hw5b, type=short, points=1, answer={0}]{3cm}{13pt}}

  \item $\ln(e^{-2}) = $ \smash{\gradedblank[id=hw5c, type=short, points=1, answer={-2}]{3cm}{13pt}}
\end{enumerate}

\item \textbf{AP-style multiple choice.} Which of the following is equivalent to $\log_2(8x^4)$?

\vspace{6pt}
\begin{enumerate}[leftmargin=2em, itemsep=4pt, label=(\Alph*)]
  \item $3 + 4\log_2(x)$
  \item $3 \cdot 4 \cdot \log_2(x)$
  \item $3 + \log_2(x^4)$
  \item $12\log_2(x)$ \hfill\textit{Letter:}~\gradedblank[id=hw6, type=choice, points=1, answer=A]{0.42in}{13pt}
\end{enumerate}

\end{enumerate}

\vspace{0.12in}

\begin{practicebox}
\textbf{Extension (optional).} Prove the quotient rule:
$\log_b\!\left(\dfrac{x}{y}\right) = \log_b(x) - \log_b(y)$.

\textit{Hint:} Write $x/y = x \cdot y^{-1}$, then apply the product property and the fact
that $b^{-n} = 1/b^n$ implies $\log_b(y^{-1}) = -\log_b(y)$.

\writelines{4}
\end{practicebox}

\end{document}
